% Options for packages loaded elsewhere
\PassOptionsToPackage{unicode}{hyperref}
\PassOptionsToPackage{hyphens}{url}
\documentclass[
]{article}
\usepackage{xcolor}
\usepackage[margin=1in]{geometry}
\usepackage{amsmath,amssymb}
\setcounter{secnumdepth}{-\maxdimen} % remove section numbering
\usepackage{iftex}
\ifPDFTeX
  \usepackage[T1]{fontenc}
  \usepackage[utf8]{inputenc}
  \usepackage{textcomp} % provide euro and other symbols
\else % if luatex or xetex
  \usepackage{unicode-math} % this also loads fontspec
  \defaultfontfeatures{Scale=MatchLowercase}
  \defaultfontfeatures[\rmfamily]{Ligatures=TeX,Scale=1}
\fi
\usepackage{lmodern}
\ifPDFTeX\else
  % xetex/luatex font selection
\fi
% Use upquote if available, for straight quotes in verbatim environments
\IfFileExists{upquote.sty}{\usepackage{upquote}}{}
\IfFileExists{microtype.sty}{% use microtype if available
  \usepackage[]{microtype}
  \UseMicrotypeSet[protrusion]{basicmath} % disable protrusion for tt fonts
}{}
\makeatletter
\@ifundefined{KOMAClassName}{% if non-KOMA class
  \IfFileExists{parskip.sty}{%
    \usepackage{parskip}
  }{% else
    \setlength{\parindent}{0pt}
    \setlength{\parskip}{6pt plus 2pt minus 1pt}}
}{% if KOMA class
  \KOMAoptions{parskip=half}}
\makeatother
\usepackage{color}
\usepackage{fancyvrb}
\newcommand{\VerbBar}{|}
\newcommand{\VERB}{\Verb[commandchars=\\\{\}]}
\DefineVerbatimEnvironment{Highlighting}{Verbatim}{commandchars=\\\{\}}
% Add ',fontsize=\small' for more characters per line
\usepackage{framed}
\definecolor{shadecolor}{RGB}{248,248,248}
\newenvironment{Shaded}{\begin{snugshade}}{\end{snugshade}}
\newcommand{\AlertTok}[1]{\textcolor[rgb]{0.94,0.16,0.16}{#1}}
\newcommand{\AnnotationTok}[1]{\textcolor[rgb]{0.56,0.35,0.01}{\textbf{\textit{#1}}}}
\newcommand{\AttributeTok}[1]{\textcolor[rgb]{0.13,0.29,0.53}{#1}}
\newcommand{\BaseNTok}[1]{\textcolor[rgb]{0.00,0.00,0.81}{#1}}
\newcommand{\BuiltInTok}[1]{#1}
\newcommand{\CharTok}[1]{\textcolor[rgb]{0.31,0.60,0.02}{#1}}
\newcommand{\CommentTok}[1]{\textcolor[rgb]{0.56,0.35,0.01}{\textit{#1}}}
\newcommand{\CommentVarTok}[1]{\textcolor[rgb]{0.56,0.35,0.01}{\textbf{\textit{#1}}}}
\newcommand{\ConstantTok}[1]{\textcolor[rgb]{0.56,0.35,0.01}{#1}}
\newcommand{\ControlFlowTok}[1]{\textcolor[rgb]{0.13,0.29,0.53}{\textbf{#1}}}
\newcommand{\DataTypeTok}[1]{\textcolor[rgb]{0.13,0.29,0.53}{#1}}
\newcommand{\DecValTok}[1]{\textcolor[rgb]{0.00,0.00,0.81}{#1}}
\newcommand{\DocumentationTok}[1]{\textcolor[rgb]{0.56,0.35,0.01}{\textbf{\textit{#1}}}}
\newcommand{\ErrorTok}[1]{\textcolor[rgb]{0.64,0.00,0.00}{\textbf{#1}}}
\newcommand{\ExtensionTok}[1]{#1}
\newcommand{\FloatTok}[1]{\textcolor[rgb]{0.00,0.00,0.81}{#1}}
\newcommand{\FunctionTok}[1]{\textcolor[rgb]{0.13,0.29,0.53}{\textbf{#1}}}
\newcommand{\ImportTok}[1]{#1}
\newcommand{\InformationTok}[1]{\textcolor[rgb]{0.56,0.35,0.01}{\textbf{\textit{#1}}}}
\newcommand{\KeywordTok}[1]{\textcolor[rgb]{0.13,0.29,0.53}{\textbf{#1}}}
\newcommand{\NormalTok}[1]{#1}
\newcommand{\OperatorTok}[1]{\textcolor[rgb]{0.81,0.36,0.00}{\textbf{#1}}}
\newcommand{\OtherTok}[1]{\textcolor[rgb]{0.56,0.35,0.01}{#1}}
\newcommand{\PreprocessorTok}[1]{\textcolor[rgb]{0.56,0.35,0.01}{\textit{#1}}}
\newcommand{\RegionMarkerTok}[1]{#1}
\newcommand{\SpecialCharTok}[1]{\textcolor[rgb]{0.81,0.36,0.00}{\textbf{#1}}}
\newcommand{\SpecialStringTok}[1]{\textcolor[rgb]{0.31,0.60,0.02}{#1}}
\newcommand{\StringTok}[1]{\textcolor[rgb]{0.31,0.60,0.02}{#1}}
\newcommand{\VariableTok}[1]{\textcolor[rgb]{0.00,0.00,0.00}{#1}}
\newcommand{\VerbatimStringTok}[1]{\textcolor[rgb]{0.31,0.60,0.02}{#1}}
\newcommand{\WarningTok}[1]{\textcolor[rgb]{0.56,0.35,0.01}{\textbf{\textit{#1}}}}
\usepackage{graphicx}
\makeatletter
\newsavebox\pandoc@box
\newcommand*\pandocbounded[1]{% scales image to fit in text height/width
  \sbox\pandoc@box{#1}%
  \Gscale@div\@tempa{\textheight}{\dimexpr\ht\pandoc@box+\dp\pandoc@box\relax}%
  \Gscale@div\@tempb{\linewidth}{\wd\pandoc@box}%
  \ifdim\@tempb\p@<\@tempa\p@\let\@tempa\@tempb\fi% select the smaller of both
  \ifdim\@tempa\p@<\p@\scalebox{\@tempa}{\usebox\pandoc@box}%
  \else\usebox{\pandoc@box}%
  \fi%
}
% Set default figure placement to htbp
\def\fps@figure{htbp}
\makeatother
\setlength{\emergencystretch}{3em} % prevent overfull lines
\providecommand{\tightlist}{%
  \setlength{\itemsep}{0pt}\setlength{\parskip}{0pt}}
\usepackage{bookmark}
\IfFileExists{xurl.sty}{\usepackage{xurl}}{} % add URL line breaks if available
\urlstyle{same}
\hypersetup{
  pdftitle={Extrusión de filamentos de PLA para fabricación aditiva},
  pdfauthor={Andres Stiven Carreño Jerez},
  hidelinks,
  pdfcreator={LaTeX via pandoc}}

\title{Extrusión de filamentos de PLA para fabricación aditiva}
\author{Andres Stiven Carreño Jerez}
\date{}

\begin{document}
\maketitle

\section{Construção de planejamento
fracionário}\label{construuxe7uxe3o-de-planejamento-fracionuxe1rio}

\begin{Shaded}
\begin{Highlighting}[]
\FunctionTok{library}\NormalTok{(FrF2)}
\end{Highlighting}
\end{Shaded}

\begin{verbatim}
## Cargando paquete requerido: DoE.base
\end{verbatim}

\begin{verbatim}
## Cargando paquete requerido: grid
\end{verbatim}

\begin{verbatim}
## Cargando paquete requerido: conf.design
\end{verbatim}

\begin{verbatim}
## Registered S3 method overwritten by 'DoE.base':
##   method           from       
##   factorize.factor conf.design
\end{verbatim}

\begin{verbatim}
## 
## Adjuntando el paquete: 'DoE.base'
\end{verbatim}

\begin{verbatim}
## The following objects are masked from 'package:stats':
## 
##     aov, lm
\end{verbatim}

\begin{verbatim}
## The following object is masked from 'package:graphics':
## 
##     plot.design
\end{verbatim}

\begin{verbatim}
## The following object is masked from 'package:base':
## 
##     lengths
\end{verbatim}

\begin{Shaded}
\begin{Highlighting}[]
\NormalTok{fatorial1}\OtherTok{=}\FunctionTok{FrF2}\NormalTok{(}\AttributeTok{nruns =} \DecValTok{8}\NormalTok{, }\AttributeTok{nfactors =} \DecValTok{3}\NormalTok{ , }\AttributeTok{factor.names =} \FunctionTok{list}\NormalTok{(}\AttributeTok{Temp =} \FunctionTok{c}\NormalTok{(}\DecValTok{180}\NormalTok{, }\DecValTok{200}\NormalTok{), }\AttributeTok{Rot =} \FunctionTok{c}\NormalTok{(}\DecValTok{30}\NormalTok{, }\DecValTok{60}\NormalTok{), }\AttributeTok{Vel =} \FunctionTok{c}\NormalTok{(}\FloatTok{1.0}\NormalTok{, }\FloatTok{2.0}\NormalTok{)),}
                 \AttributeTok{randomize =} \ConstantTok{FALSE}\NormalTok{)}
\end{Highlighting}
\end{Shaded}

\begin{verbatim}
## creating full factorial with 8 runs ...
\end{verbatim}

\section{Digite a resposta
(diâmetro)}\label{digite-a-resposta-diuxe2metro}

\begin{Shaded}
\begin{Highlighting}[]
\NormalTok{diam}\OtherTok{=} \FunctionTok{c}\NormalTok{(}\FloatTok{1.92}\NormalTok{, }\FloatTok{1.86}\NormalTok{, }\FloatTok{2.05}\NormalTok{, }\FloatTok{1.98}\NormalTok{, }\FloatTok{1.70}\NormalTok{, }\FloatTok{1.65}\NormalTok{, }\FloatTok{1.82}\NormalTok{, }\FloatTok{1.76}\NormalTok{)}
\NormalTok{fatorial2}\OtherTok{=}\FunctionTok{add.response}\NormalTok{(fatorial1,diam)}
\end{Highlighting}
\end{Shaded}

\begin{Shaded}
\begin{Highlighting}[]
\FunctionTok{print}\NormalTok{(fatorial2)}
\end{Highlighting}
\end{Shaded}

\begin{verbatim}
##   Temp Rot Vel diam
## 1  180  30   1 1.92
## 2  200  30   1 1.86
## 3  180  60   1 2.05
## 4  200  60   1 1.98
## 5  180  30   2 1.70
## 6  200  30   2 1.65
## 7  180  60   2 1.82
## 8  200  60   2 1.76
## class=design, type= full factorial
\end{verbatim}

\section{Ajuste um modelo onde o diâmetro seja uma função de A, B, C e
todas as suas
interações.}\label{ajuste-um-modelo-onde-o-diuxe2metro-seja-uma-funuxe7uxe3o-de-a-b-c-e-todas-as-suas-interauxe7uxf5es.}

\begin{Shaded}
\begin{Highlighting}[]
\NormalTok{mod\_compl}\OtherTok{=}\FunctionTok{lm}\NormalTok{(diam }\SpecialCharTok{\textasciitilde{}}\NormalTok{ Temp }\SpecialCharTok{*}\NormalTok{ Rot }\SpecialCharTok{*}\NormalTok{ Vel, }\AttributeTok{data =}\NormalTok{ fatorial1)}
\FunctionTok{summary}\NormalTok{(mod\_compl)}
\end{Highlighting}
\end{Shaded}

\begin{verbatim}
## 
## Call:
## lm.default(formula = diam ~ Temp * Rot * Vel, data = fatorial1)
## 
## Residuals:
## ALL 8 residuals are 0: no residual degrees of freedom!
## 
## Coefficients:
##                   Estimate Std. Error t value Pr(>|t|)
## (Intercept)      1.842e+00        NaN     NaN      NaN
## Temp1           -3.000e-02        NaN     NaN      NaN
## Rot1             6.000e-02        NaN     NaN      NaN
## Vel1            -1.100e-01        NaN     NaN      NaN
## Temp1:Rot1      -2.500e-03        NaN     NaN      NaN
## Temp1:Vel1       2.500e-03        NaN     NaN      NaN
## Rot1:Vel1       -2.500e-03        NaN     NaN      NaN
## Temp1:Rot1:Vel1  2.208e-17        NaN     NaN      NaN
## 
## Residual standard error: NaN on 0 degrees of freedom
## Multiple R-squared:      1,  Adjusted R-squared:    NaN 
## F-statistic:   NaN on 7 and 0 DF,  p-value: NA
\end{verbatim}

\section{↑↑↑ Modelo saturado, e sempre ocorre em um desenho 2³ sem
réplicas.}\label{modelo-saturado-e-sempre-ocorre-em-um-desenho-2uxb3-sem-ruxe9plicas.}

\section{↓↓↓ Gráfico (meio-normal) para identificar efeitos
importantes.}\label{gruxe1fico-meio-normal-para-identificar-efeitos-importantes.}

\begin{Shaded}
\begin{Highlighting}[]
\FunctionTok{DanielPlot}\NormalTok{(mod\_compl)}
\end{Highlighting}
\end{Shaded}

\pandocbounded{\includegraphics[keepaspectratio]{1eje_files/1eje_files/figure-latex/unnamed-chunk-5-1.pdf}}

\begin{Shaded}
\begin{Highlighting}[]
\NormalTok{Efeitos}\OtherTok{=}\DecValTok{2}\SpecialCharTok{*}\NormalTok{mod\_compl}\SpecialCharTok{$}\NormalTok{coefficients}

\NormalTok{Efeitos}
\end{Highlighting}
\end{Shaded}

\begin{verbatim}
##     (Intercept)           Temp1            Rot1            Vel1      Temp1:Rot1 
##    3.685000e+00   -6.000000e-02    1.200000e-01   -2.200000e-01   -5.000000e-03 
##      Temp1:Vel1       Rot1:Vel1 Temp1:Rot1:Vel1 
##    5.000000e-03   -5.000000e-03    4.415885e-17
\end{verbatim}

\section{↓↓↓ Uma apresentação
melhor}\label{uma-apresentauxe7uxe3o-melhor}

\begin{Shaded}
\begin{Highlighting}[]
\NormalTok{Efeitos}\OtherTok{=}\DecValTok{2}\SpecialCharTok{*}\NormalTok{mod\_compl}\SpecialCharTok{$}\NormalTok{coefficients}

\FunctionTok{data.frame}\NormalTok{(  }\AttributeTok{Efecto =} \FunctionTok{c}\NormalTok{(}\StringTok{"Media"}\NormalTok{, }\StringTok{"A(Temp)"}\NormalTok{, }\StringTok{"B(Rot)"}\NormalTok{, }\StringTok{"C(Vel)"}\NormalTok{,}\StringTok{"AB"}\NormalTok{, }\StringTok{"AC"}\NormalTok{, }\StringTok{"BC"}\NormalTok{, }\StringTok{"ABC"}\NormalTok{), }\AttributeTok{Valor  =} \FunctionTok{round}\NormalTok{(Efeitos, }\DecValTok{5}\NormalTok{))}
\end{Highlighting}
\end{Shaded}

\begin{verbatim}
##                  Efecto  Valor
## (Intercept)       Media  3.685
## Temp1           A(Temp) -0.060
## Rot1             B(Rot)  0.120
## Vel1             C(Vel) -0.220
## Temp1:Rot1           AB -0.005
## Temp1:Vel1           AC  0.005
## Rot1:Vel1            BC -0.005
## Temp1:Rot1:Vel1     ABC  0.000
\end{verbatim}

\section{Efeitos Principais e
Interações}\label{efeitos-principais-e-interauxe7uxf5es}

\section{Gráfico de Efeitos
Principais}\label{gruxe1fico-de-efeitos-principais}

\begin{Shaded}
\begin{Highlighting}[]
\FunctionTok{MEPlot}\NormalTok{(fatorial2)}
\end{Highlighting}
\end{Shaded}

\pandocbounded{\includegraphics[keepaspectratio]{1eje_files/1eje_files/figure-latex/unnamed-chunk-8-1.pdf}}
\# Gráfico de interação

\begin{Shaded}
\begin{Highlighting}[]
\FunctionTok{IAPlot}\NormalTok{(fatorial2)}
\end{Highlighting}
\end{Shaded}

\pandocbounded{\includegraphics[keepaspectratio]{1eje_files/1eje_files/figure-latex/unnamed-chunk-9-1.pdf}}

\section{¿Qué ensayo produce diámetro más cercano a 1,75 mm?
---}\label{quuxe9-ensayo-produce-diuxe1metro-muxe1s-cercano-a-175-mm}

\begin{Shaded}
\begin{Highlighting}[]
\NormalTok{fatorial2}\SpecialCharTok{$}\NormalTok{error\_nominal}\OtherTok{=}\FunctionTok{abs}\NormalTok{(fatorial2}\SpecialCharTok{$}\NormalTok{diam }\SpecialCharTok{{-}} \FloatTok{1.75}\NormalTok{)}
\FunctionTok{cat}\NormalTok{(}\StringTok{"}\SpecialCharTok{\textbackslash{}n}\StringTok{{-}{-}{-} Ensayo más cercano a 1,75 mm {-}{-}{-}}\SpecialCharTok{\textbackslash{}n}\StringTok{"}\NormalTok{)}
\end{Highlighting}
\end{Shaded}

\begin{verbatim}
## 
## --- Ensayo más cercano a 1,75 mm ---
\end{verbatim}

\begin{Shaded}
\begin{Highlighting}[]
\FunctionTok{print}\NormalTok{(fatorial2[}\FunctionTok{order}\NormalTok{(fatorial2}\SpecialCharTok{$}\NormalTok{error\_nominal), }\FunctionTok{c}\NormalTok{(}\StringTok{"Temp"}\NormalTok{,}\StringTok{"Rot"}\NormalTok{,}\StringTok{"Vel"}\NormalTok{,}\StringTok{"diam"}\NormalTok{,}\StringTok{"error\_nominal"}\NormalTok{)])}
\end{Highlighting}
\end{Shaded}

\begin{verbatim}
##   Temp Rot Vel diam error_nominal
## 8  200  60   2 1.76          0.01
## 5  180  30   2 1.70          0.05
## 7  180  60   2 1.82          0.07
## 6  200  30   2 1.65          0.10
## 2  200  30   1 1.86          0.11
## 1  180  30   1 1.92          0.17
## 4  200  60   1 1.98          0.23
## 3  180  60   1 2.05          0.30
\end{verbatim}

El factor más importante fue: C -- Velocidad de arrastre (efecto =
−0,220 mm)

\section{--- Su efecto NEGATIVO indica que pasar del nivel
bajo}\label{su-efecto-negativo-indica-que-pasar-del-nivel-bajo}

al nivel alto REDUCE significativamente el diámetro. En este caso: -
Nivel bajo = 1,0 m/min - Nivel alto = 2,0 m/min → Mayor velocidad de
arrastre produce filamentos más delgados.

\section{--- Los demás factores
presentaron:}\label{los-demuxe1s-factores-presentaron}

\begin{itemize}
\tightlist
\item
  B (Rotación): efecto moderado positivo (+0,120 mm) Mayor rotación →
  más material → filamento más grueso
\item
  A (Temperatura): efecto pequeño negativo (−0,060 mm) Mayor temperatura
  → menor viscosidad → filamento más delgado
\item
  AB, AC, BC, ABC: efectos despreciables (≤ 0,005 mm) Los factores
  actúan de forma independiente entre sí
\end{itemize}

\section{--- La mejor condición experimental para
alcanzar}\label{la-mejor-condiciuxf3n-experimental-para-alcanzar}

el diámetro nominal de 1,75 mm fue el ENSAYO 8: - Temperatura : 200 °C
(nivel alto) - Rotación : 60 rpm (nivel alto) - Vel. arrastre: 2,0 m/min
(nivel alto) - Diámetro obtenido : 1,76 mm - Error al nominal : 0,01 mm
← menor de todos

\section{--- Esta condición también
presentó:}\label{esta-condiciuxf3n-tambiuxe9n-presentuxf3}

\begin{itemize}
\tightlist
\item
  Menor desviación estándar del diámetro (DP = 0,040 mm)
\item
  Menor ovalización (2,1\%) → Es la condición más precisa Y más exacta
  del experimento.
\end{itemize}

\end{document}
